\documentclass[11pt]{article}
\usepackage{amsmath,amssymb}
\usepackage[margin=0.94in]{geometry}
\setlength{\parindent}{0pt}
\setlength{\parskip}{5pt}
\setlength{\emergencystretch}{2em}
\newcommand{\R}{\mathbb R}
\newcommand{\F}{\mathcal F}
\title{Two symmetry-breaking bifurcations of planar\newline Newtonian six-body central configurations}
\author{Chatgpt}
\date{Research manuscript, 27 November 2026}
\begin{document}
\maketitle
\begin{abstract}
We certify two pitchforks on the mass line $(1,1,1,1,1,t)$ at
$t\simeq0.089353633682323$ and $t\simeq0.175165824529870$.
For each of the two shapes, the local count of labelled, oriented central
configurations changes from three to one as $t$ increases. Splitting the
masses of a reflected pair provides a second physical parameter that gives a
versal cusp at each point. Near either shape, every sufficiently close
six-tuple of independently perturbed masses has at most three solutions.
An integer interval calculation isolates the two degeneracies and bounds the
derivatives of the reduced equation. All counts are local in configuration
space.
\end{abstract}

\section{Statements and configuration geometry}
A central configuration is a collision-free set of particle positions whose
Newtonian accelerations are proportional to their displacements towards
the centre of mass. Smale's sixth problem asks for finiteness up to
similarity for each fixed choice of positive masses. Even complete counts
at individual mass vectors can miss changes as mass ratios vary. At the
six-equal-mass vector, Moczurad and Zgliczy\'nski found every planar shape
to have a line of reflection symmetry \cite{mz}. Dias and Pan \cite{dp}
proved generic finiteness among six-body shapes with precisely four axial
particles, which they call cross central configurations. In the four-body
problem Rusu and Santoprete \cite{rs} used interval arithmetic and
equivariant bifurcation theory to classify mass bifurcations. We study two specific
transitions on a one-parameter family of unequal masses, and the effect of
varying these masses independently.

\textbf{Theorem.}
Let the first five masses be $1$ and the sixth mass be $t>0$.
Write
\[
t_1^0=0.089353633682322681683088206548106878,\quad
t_2^0=0.17516582452986998163624081143199902.
\]
For $i=1,2$ there is a collision-free central configuration $Q_i$ at
$|t_i-t_i^0|\le10^{-15}$. Each $Q_i$ has four bodies on a reflection axis
and a reflected pair off that axis. In a neighbourhood of $Q_i$ modulo
direct similarity, for $t$ sufficiently close to $t_i$ there are exactly three
labelled configurations when $t<t_i$ and exactly one when $t\ge t_i$.
The two additional configurations break the reflection symmetry.
At any sufficiently close vector of six independently varying positive masses,
there are at most three configurations in this neighbourhood. Varying $t$ and
splitting the masses of the off-axis pair as $(1+\delta,1-\delta)$ gives a
two-parameter cusp unfolding.

The two tips are on geometrically different configurations. Among the four
axial particles, the light particle is second from the left at the first tip
and leftmost at the second tip in the chart below. Labels and mirror images
are retained in the count. No enumeration of shapes away from these two
neighbourhoods enters the argument.

\newpage
\section{Equations, rotational slice and reflection}
Divide all masses by five and use length units in which the acceleration
multiplier is one. Number the particles $0,\ldots,5$. For positive masses
$m_i$, set
\[
 F_i(q,m)=q_i-\sum_{j\ne i}m_j
 \frac{q_i-q_j}{|q_i-q_j|^3}.
\]
If $F=0$, summing the mass-weighted equations gives zero centre of mass.
In a chart with $q_{3y}=0$ and $q_{3x}>0$, drop the equation $F_{3y}=0$.
The identity $\sum_i m_i q_i\wedge F_i=0$ recovers that equation from all
the others. Let $\F$ denote this eleven-dimensional residual in the
eleven-dimensional chart. It also applies when the masses vary independently.

Reflection across the horizontal axis combined with the transposition $(1\ 4)$
defines an involution $S$. At equal masses for particles 1 and 4, $\F$ is
$S$-equivariant. In both the position slice and residual space the even and odd
subspaces have dimensions six and five. In zero-based coordinates the
embeddings used by the certificate are
\begin{align*}
 P(p)&=((p_0,0),(p_4,p_5),(p_1,0),(p_2,0),(p_4,-p_5),(p_3,0)),\\
 V(v)&=((0,v_0),(v_3,v_4),(0,v_1),(0,0),(-v_3,v_4),(0,v_2)).
\end{align*}
For the residual, the displayed slot $3y$ is simply absent. Denote
the Jacobian blocks at an even point by $J_+$ and $J_-$. A tip is located
by solving the twelve equations
\[
 \F_+(P(p),t)=0,\qquad J_-(p,t)v=0,\qquad v_k=1                 \tag{1}
\]
for $(p,t,v)$. For later use define the bordered odd matrix
\[
 M=\begin{pmatrix}J_-&e_b\\ e_k^{\mathsf T}&0\end{pmatrix}.
\]
The table gives rounded coordinates for identification. The JSON certificates
supply the unrounded box centres.
\[
\begin{array}{c|rrrrrr|cc}
 &p_0&p_1&p_2&p_3&p_4&p_5&k&b\\ \hline
 1&-.74733&.08957&.77496&-.33508&-.04363&-.72946&0&3\\
 2&-.62719&.09629&.78201&-1.13823&-.02586&-.73844&3&3
\end{array}
\]
At both tips the boxes certify positive pairwise distances, $q_{3x}>0$ and
an anchor $q_3$ with squared radius distinct from those of the other heavy
particles. In a small neighbourhood any mass-preserving reflection must
fix $q_3$ and the centre of mass, so its axis must be horizontal. The only
possible reflection permutation there is $(1\ 4)$, by the separation of
the particles. Breaking $S$ therefore removes all physical reflection axes
along the equal-heavy-mass line.

\newpage
\section{Range elimination and the local counts}
The checker proves that $J_+$ and $M$ are invertible at each tip. Thus $J_-$
has a one-dimensional kernel and cokernel. Define $\rho$ on odd residuals
by the last component of $M^{-1}(y,0)$, and extend it by zero on even
residuals. Then $\rho J_-=0$ and $\rho e_b=1$. Let $p(t)$ be the symmetric
solution continued from the tip by the implicit function theorem. Put
$s=e_k^{\mathsf T}x_-$, where $x_-$ is the odd position coordinate;
let $v$ also denote its physical embedding $V(v)$. In the coordinates
centred on $P(p(t))$, solve the projected equations for the ten
complementary directions by the analytic implicit function theorem. The
projection subtracts $e_b\rho$. The remaining scalar residual on this
solution manifold is $g(s,t,\delta)=\rho\F$. At $\delta=0$ uniqueness of
the range solution makes $g$ odd in $s$.

Let $w= dP(p(t))/dt$ and $u$ be the second $s$-derivative of the position
on the range manifold at the tip. Both are even there. They are found
from $J_+w=-\F_t$ and $J_+u=-D^2\F[v,v]$ in even coordinates. At the
tip the checker evaluates the following scalars in the indicated
normalization:
\begin{align*}
 A=g_{st}&=\rho\big((\partial_tD\F)v+D^2\F[v,w]\big),\\
 C=g_{sss}&=\rho\big(D^3\F[v,v,v]+3D^2\F[v,u]\big),\\
 B=g_{\delta}&=\rho\F_{\delta}.
\end{align*}
All omitted linear range-correction terms are annihilated by $\rho$.
The following coarsened rational enclosures contain the signed bounds
printed by the checker.
\[
\begin{array}{c|ccc}
 & A&C&B\\ \hline
1 &[1.56069,1.56070]&[11.08269,11.08270]&[-.14608,-.14607]\\
2 &[-1.15635,-1.15634]&[-18.27376,-18.27375]&[-.14127,-.14126]
\end{array}
\]
On the symmetric mass line write $g=s h(s^2,t)$. At $(0,t_i)$ one has
$h=0$, $h_t=A$, $h_z=C/6$. The nontrivial roots satisfy
$t=\tau(s^2)$ with $\tau'(0)=-C/(6A)<0$.
This proves the local 3-to-1 count, including the isolated triple root
at the critical mass.

For all six mass parameters the same projected equations give a scalar
real-analytic residual with an order-three zero in $s$ at the base point.
Weierstrass preparation \cite{lebl} (complexify the real-analytic germ and
use uniqueness for reality) factors it as a nonvanishing unit times
$s^3+c_2s^2+c_1s+c_0$, with analytic mass-dependent coefficients vanishing
at the tip. This gives the local upper bound three. On the two-parameter
slice $(t,\delta)$, translate $s$ to depress the cubic to $z^3+a z+b$.
Writing $U$ for the unit and using $g_t=0$ at the tip gives
$a_t=A/U$, $b_t=0$, $b_\delta=B/U$ there, where $U=C/6$.
Consequently $(a,b)$ are local mass coordinates. There are three distinct
local real shapes when $4a^3+27b^2<0$, one when this expression is
positive, two on its smooth zero set and one at the tip.

\newpage
\section{Reproducing the interval proof}
Run \texttt{python3 verify.py} with \texttt{inner.json} and \texttt{data.json}
in the same directory. Each file contains the centre $c$ of a
twelve-dimensional cube, the two indices $k,b$, and a rational
preconditioner $R$. All entries are finite decimal strings. The program
uses standard-library integers, fractions and integer square roots. It
encloses real arithmetic on the rational lattice $10^{-42}\mathbb Z$ with
outward rounding. Candidate discovery and floating-point inverses do not enter
the verification.

Write $G$ for the left side of (1), subtracting 1 in its last component.
Let $X$ be the closed cube about $c$ of radius $r=10^{-15}$ in the
supremum norm. The checker encloses $DG(X)$, and bounds the infinity
matrix norm by summing absolute entry bounds in each row. In both cubes,
\[
 \|R G(c)\|_\infty<10^{-25},\qquad
 \sup_{x\in X}\|I-RDG(x)\|_\infty<5\cdot10^{-12}.
\]
Thus $x\mapsto x-RG(x)$ maps $X$ strictly into itself and is a contraction.
The norm bound implies that $RDG(c)$, hence the square matrix $R$, is
invertible. The fixed point gives a unique solution of (1) in $X$. For
the matrices $J_+$ and $M$, interval Gaussian elimination verifies
pivots separated from zero and encloses all required solves over $X$.
Bounds on the coefficients and on the locations of the tips follow.

The fold curves can also be read in the physical mass coordinates.
Solving $g=g_s=0$ for $(t-t_i,\delta)$ as functions of $s$ is possible
since $AB\ne0$. To leading order this gives
$t-t_i\sim-C s^2/(2A)$ and $\delta\sim C s^3/(3B)$.
Hence the two fold curves have
$\delta\sim\pm K(t_i-t)^{3/2}$, with
$K=\sqrt{8A^3/(9B^2C)}$. The corresponding half-width coefficients are
$K\simeq3.780049$ and $K\simeq1.941401$.

The rational enclosures of the squared radii and pairwise distances also
verify the geometric hypotheses of the chart. This proof establishes no
quantitative size for the continuation neighbourhood in the independent
mass coordinates. To establish the full six-body finiteness conjecture
one would have to control every configuration at every positive mass
vector, including any singular configurations not covered by these cubes.

\begin{thebibliography}{9}
\bibitem{mz}
M. Moczurad and P. Zgliczy\'nski,
Central configurations in planar $n$-body problem with equal masses for
$n=5,6,7$. \emph{Celestial Mechanics and Dynamical Astronomy}
\textbf{131}, article 46 (2019). DOI: 10.1007/s10569-019-9920-6.
\bibitem{dp} T. Dias and B.-Y. Pan,
Generic finiteness for a class of symmetric planar central configurations
of the six-body problem and the six-vortex problem.
\emph{Journal of Dynamics and Differential Equations}
\textbf{32}, 1579--1602 (2020). DOI: 10.1007/s10884-019-09765-y.
\bibitem{rs} D. Rusu and M. Santoprete,
Bifurcations of central configurations in the four-body problem with some equal masses.
\emph{SIAM Journal on Applied Dynamical Systems} (2016).
DOI: 10.1137/151002484.
\bibitem{lebl} J. Lebl, \emph{Tasty Bits of Several Complex Variables},
Section 6.2, Weierstrass preparation and division theorems.
\end{thebibliography}
\medskip\noindent Original manuscript and source: CC BY 4.0. Verification code: MIT.
\end{document}
